\section*{第 \ref{ch:b1:rate-laws} 章——化学动力学：速率方程}

\begin{solution}{exo:b1:rate-laws:1}
$v = \frac14 \times 2.0 \times 10^{-4} = \qty{5.0e-5}{mol.L^{-1}.s^{-1}}$；
氧气以相同速率生成，为 \qty{5.0e-5}{}；\ce{N2O5} 的\omterm{def:b1:rate-laws:rate}{消耗速率}为
$2v = \qty{1.0e-4}{mol.L^{-1}.s^{-1}}$。
\end{solution}

\begin{solution}{exo:b1:rate-laws:2}
零级：\unit{mol.L^{-1}.s^{-1}}；一级：\unit{s^{-1}}；二级：
\unit{L.mol^{-1}.s^{-1}}；三级：\unit{L^2.mol^{-2}.s^{-1}}。
\end{solution}

\begin{solution}{exo:b1:rate-laws:3}
$t_{1/2} = \ln2/k = 0.693/(3.0 \times 10^{-3}) = \qty{231}{s}$。经过
\qty{600}{s} 后，$\eu^{-kt} = \eu^{-1.8} = 0.165$：剩余 16.5\,\%。
\end{solution}

\begin{solution}{exo:b1:rate-laws:4}
\omterm{def:b1:rate-laws:half-life}{半衰期}与 $1/a_0$ 成正比，故为二级反应。若 $a_0$ 减半时\omterm{def:b1:rate-laws:half-life}{半衰期}也减半，
则 $t_{1/2} \propto a_0$：为零级反应。
\end{solution}

\begin{solution}{exo:b1:rate-laws:5}
$[\ce{A}]$ 与 $t$ 不呈线性关系（减少量依次为 0.0259、0.0192、0.0142、0.0106）。
$\ln[\ce{A}] = -2.303$、$-2.602$、$-2.902$、$-3.202$、$-3.503$：每隔
\qty{10}{min} 下降 0.300，呈直线关系。为一级反应，$k =
\qty{0.0300}{min^{-1}}$。
\end{solution}

\begin{solution}{exo:b1:rate-laws:6}
将 $[\ce{A}]_0$ 加倍，$v_0$ 变为 $4 = 2^2$ 倍：对 \ce{A} 为二级。
将 $[\ce{B}]_0$ 加倍，速率变为 2 倍：对 \ce{B} 为一级。$v =
k[\ce{A}]^2[\ce{B}]$，$k = 2.0 \times 10^{-6}/(0.010^2 \times 0.010) =
\qty{2.0}{L^2.mol^{-2}.s^{-1}}$。
\end{solution}

\begin{solution}{exo:b1:rate-laws:7}
\ce{B} 大大过量：其浓度保持为
\qty{0.50}{mol/L}，且 $v = k'[\ce{A}]$，其中 $k' = k[\ce{B}]_0$。因此
$k = 0.012/0.50 = \qty{0.024}{L.mol^{-1}.s^{-1}}$。
\end{solution}

\begin{solution}{exo:b1:rate-laws:8}
$E_a = R\ln5/(1/300 - 1/320) = 8.314 \times 1.609/(2.083 \times 10^{-4}) =
\qty{64.2}{kJ/mol}$。在 \qty{310}{K} 时：$\ln k = \ln(1.0 \times 10^{-3}) +
\frac{E_a}{R}\big(\frac{1}{300} - \frac{1}{310}\big) = -6.908 + 0.831$，
$k = \qty{2.3e-3}{s^{-1}}$。
\end{solution}

\begin{solution}{exo:b1:rate-laws:9}
此时 $a = 2$：$1/[\ce{A}] = 1/a_0 + 2kt$，$t_{1/2} = 1/(2k a_0) = 1/(2
\times 0.50 \times 0.020) = \qty{50}{s}$。在 \qty{100}{s} 时：$1/[\ce{A}] = 50
+ 100 = \qty{150}{L/mol}$，$[\ce{A}] = \qty{6.7e-3}{mol/L}$。
\end{solution}

\begin{solution}{exo:b1:rate-laws:10}
$[\ce{A}] = a_0(1 - f) = a_0\eu^{-akt_f}$ 给出 $t_f = -\ln(1-f)/(ak)$，
以\omterm{def:b1:rate-laws:half-life}{半衰期}计，$t_f/t_{1/2} = -\ln(1-f)/\ln2$：完成 50\,\% 时为 1，75\,\% 时为 2，
90\,\% 时为 3.32，99.9\,\% 时为 9.97；约十个\omterm{def:b1:rate-laws:half-life}{半衰期}后反应
可以“完成”。
\end{solution}

\begin{solution}{exo:b1:rate-laws:11}
$m(t) = 10 - 0.40\,t$（\unit{mg}，\unit{h}）：为零级反应。\omterm{def:b1:rate-laws:half-life}{半衰期}
$10/(2 \times 0.40) = \qty{12.5}{h}$；\qty{25}{h} 后药量耗尽。每小时释放的
剂量与剩余量无关，即持续稳定供药，
这正是贴片的用途。
\end{solution}

\begin{solution}{exo:b1:rate-laws:12}
$E_a = R\ln2/(1/298.15 - 1/308.15) = 8.314 \times 0.693/(1.088 \times
10^{-4}) = \qty{52.9}{kJ/mol}$。在 273.15 与 \qty{283.15}{K} 之间：
$\exp\big(\frac{52900}{8.314}(\frac{1}{273.15} - \frac{1}{283.15})\big) =
2.3$：\omterm{def:b1:rate-laws:activation-energy}{活化能}相同时，
温度越低，倍数因子越大。
\end{solution}

\begin{solution}{pb:b1:rate-laws:1}
\textbf{1.} 损失分别为 8.6、7.9、7.2 和 6.5 个百分点：不是常数，故不是
零级反应。
\textbf{2.} $\ln$ 值：4.605、4.515、4.425、4.335、4.245；
每隔 10 天下降 0.090，呈直线关系，为一级反应。
\textbf{3.} $k_{60} = 0.090/10 = \qty{9.0e-3}{d^{-1}}$。
\textbf{4.} $t_{1/2} = \ln2/k_{60} = \qty{77}{d}$。
\textbf{5.} $100\,\eu^{-0.54} = 58.3\,\%$。
\textbf{6.} $t_{90} = \ln(10/9)/k_{60} = 0.1054/0.0090 = \qty{11.7}{d}$。
\textbf{7.} 剩余 90\,\% 时：$\eu^{-kt_{90}} = 0.9$，故 $t_{90} = \ln(10/9)/k$。
\textbf{8.} $t_{90}/t_{1/2} = \ln(10/9)/\ln2 = 0.152$。
\textbf{9.} 对一级反应，剩余分数仅取决于 $kt$，
与初始量无关。
\textbf{10.} 在 \qty{25}{\celsius} 时，降解需要数年；加热
可加速降解（阿伦尼乌斯），使其可在数周内测得。
\textbf{11.} 对二级反应，\omterm{def:b1:rate-laws:half-life}{半衰期}为 $1/(ak\,a_0)$：药片强度加倍时，
相对降解速度加倍，保存时间变短。
\textbf{12.} $1/T$（\unit{K^{-1}}）：\num{3.1934e-3}、\num{3.0945e-3}、
\num{3.0017e-3}；$\ln k$：$-6.669$、$-5.661$、$-4.711$。
\textbf{13.} 斜率 $= (-4.711 + 6.669)/(3.0017 \times 10^{-3} - 3.1934 \times
10^{-3}) = \qty{-1.021e4}{K}$。
\textbf{14.} $E_a = 8.314 \times 1.021 \times 10^4 = \qty{85}{kJ/mol}$。
\textbf{15.} 预测 $\ln k_{50} = -4.711 - 1.021 \times 10^4 \times (3.0945 -
3.0017) \times 10^{-3} = -5.658$，实测为 $-5.661$：落在直线上。
\textbf{16.} $A = k_{60}\,\eu^{E_a/RT} = 0.0090\,\eu^{10215/333.15} =
\qty{1.9e11}{d^{-1}}$。
\textbf{17.} $\exp\big(1.021 \times 10^4(\frac{1}{298.15} -
\frac{1}{308.15})\big) = 3.0$：比规则中的因子 2 更大，因为
$E_a$ 大于 \qty{53}{kJ/mol}。
\textbf{18.} $k_{30} = k_{60}\exp\big(-1.021 \times 10^4(\frac{1}{303.15} -
\frac{1}{333.15})\big) = \qty{4.3e-4}{d^{-1}}$。
\textbf{19.} $t_{90} = 0.1054/4.33 \times 10^{-4} = \qty{243}{d}$，约 8.0
个月。
\textbf{20.} $1 - \eu^{-3 \times 0.0090} = 2.7\,\%$。
\textbf{21.} 保质期随温度迅速缩短（8 个月，
在 \qty{30}{\celsius} 时对比 \qty{25}{\celsius} 时的 14 个月）：有效期
只在低于 \qty{25}{\celsius} 时成立。
\textbf{22.} $k_{25} = k_{60}\exp\big(-1.021 \times 10^4(\frac{1}{298.15} -
\frac{1}{333.15})\big) = \qty{2.46e-4}{d^{-1}}$。
\textbf{23.} $t_{90} = 0.1054/2.46 \times 10^{-4} = \qty{428}{d}$，即
\textbf{14 个月}：印在包装盒上的日期，由
六周的测量得出。
\end{solution}
